\section{Reason for choosing the topic}

In recent years, the global fashion industry has witnessed a massive transition toward e-commerce and online clothing rental platforms, offering consumers an unprecedented variety of apparel choices \parencite{article1}. However, this rapid catalog expansion has induced a severe "information overload" problem, where users struggle to navigate extensive inventories to find items matching their personalized aesthetic preferences. Consequently, personalized Fashion Recommender Systems (FaRS) have become essential infrastructure for online retailers to enhance customer engagement and platform retention.

Unlike traditional recommender domains (e.g., movies or books), fashion recommendation presents unique domain-specific challenges. Purchasing or renting apparel is heavily driven by visual aesthetics (such as colors, patterns, and styling shapes), outfit compatibility, and subjective personal tastes. Furthermore, due to the high seasonality and short lifecycle of fashion products, transaction logs are notoriously sparse. While classical collaborative filtering (CF) algorithms rely solely on user-item interaction IDs, they suffer from severe performance degradation under physical data sparsity and fail to establish representations for newly introduced or low-interaction items.

To capture high-order collaborative relationships, Graph Neural Networks (GNNs) have emerged as the state-of-the-art paradigm in collaborative filtering. Architectures like LightGCN \parencite{he2020lightgcnsimplifyingpoweringgraph} model the recommendation task as a message-passing process over a user-item bipartite graph, capturing complex structural connections \parencite{11359174}. Nevertheless, pure GNN-based recommendation remains fundamentally ID-bound. Because they do not ingest descriptive item content, their preference propagation remains ineffective for items situated at the sparse tail of the interaction distribution.

To mitigate this structural bottleneck, multimodal recommendation has emerged as a promising research direction. By leveraging textual metadata (e.g., categories and descriptions) and visual descriptors (e.g., clothing textures and patterns), multimodal frameworks construct content-aware representations that bridge the gap between interaction signals and item semantics \parencite{10.1145/3543507.3583251}. However, dynamically fusing heterogeneous visual and textual features into GNN propagation layers while filtering out semantic noise remains an open challenge \parencite{Zhou_2023, 10.1007/978-981-95-3355-8_6}.

Motivated by these practical challenges and theoretical gaps, this study investigates, implements, and benchmarks multimodal graph collaborative filtering models on a fashion rental dataset. By systematically evaluating diverse visual encoders (MobileNetV2 vs. CLIP) and modal fusion strategies across GNN architectures, this research aims to optimize representation learning, mitigate data sparsity issues, and improve top-$N$ ranking accuracy for personalized fashion recommendation.
